\subsection{Minimal carrier theorem}
\label{sec:minimal-carrier-proof}

\begin{theorem}[Minimal carrier theorem]
Let $E=E_b\oplus E_s$ be a split carrier with traceless diagonal generator
\[
Y_{b,s}
=
\mathrm{diag}\!\left(
\underbrace{-\frac1b,\ldots,-\frac1b}_{b\ \text{times}},
\underbrace{\frac1s,\ldots,\frac1s}_{s\ \text{times}}
\right).
\]
If one simultaneously requires
\begin{equation}
\label{eq:minimal-carrier-conditions}
\dim \Lambda^{\mathrm{even}}E=16,
\qquad
\tr_E Y_{b,s}^2=\frac56,
\end{equation}
then necessarily $\{b,s\}=\{3,2\}$.
\end{theorem}

\begin{proof}
Write $g=b+s$. Since
\[
\dim \Lambda^{\mathrm{even}}E=2^{g-1},
\]
the first condition gives $g=5$. For the diagonal generator above one has
\[
\tr_E Y_{b,s}^2
=
b\left(\frac{1}{b^2}\right)+s\left(\frac{1}{s^2}\right)
=
\frac1b+\frac1s
=
\frac{b+s}{bs}
=
\frac{5}{bs}.
\]
Imposing $\tr_E Y_{b,s}^2=5/6$ yields $bs=6$. Together with $b+s=5$ this forces the unordered solution $\{b,s\}=\{3,2\}$.
\end{proof}

Hence from now on
\[
E=E_3\oplus E_2,
\qquad
P_3=\frac{\mathbf 1-\inv}{2},
\qquad
P_2=\frac{\mathbf 1+\inv}{2},
\]
and the hypercharge operator is read concretely as
\begin{equation}
\label{eq:explicit-hypercharge}
Y
=
\mathrm{diag}\!\left(-\frac13,-\frac13,-\frac13,\frac12,\frac12\right)
=
-\frac13 P_3+\frac12 P_2
=
\frac1{12}\mathbf 1+\frac5{12}\inv.
\end{equation}

\begin{lemma}[Hypercharge--seam inversion]
On the minimal carrier, the seam involution is recovered from hypercharge by
\begin{align}
\inv&=\frac{12Y-\mathbf 1}{5},
\label{eq:hypercharge-seam-inversion}\\
(12Y-\mathbf 1)^2&=25\,\mathbf 1,
\nonumber\\
6Y^2-Y-\mathbf 1&=0.
\label{eq:carrier-polynomial}
\end{align}
Consequently the carrier projectors are
\[
P_2=\frac{\mathbf 1+\inv}{2}=\frac{2(3Y+\mathbf 1)}{5},
\qquad
P_3=\frac{\mathbf 1-\inv}{2}=\frac{3(\mathbf 1-2Y)}{5}.
\]
\end{lemma}

\begin{proof}
The affine relation between $Y$ and $\inv$ is inverted algebraically. Squaring the result and using $\inv^2=\mathbf 1$ gives the quadratic carrier polynomial. The projector formulas then follow from $P_2=(\mathbf 1+\inv)/2$ and $P_3=(\mathbf 1-\inv)/2$.
\end{proof}

The hard carrier package is summarized visually in \Cref{fig:minimal-carrier-glance}. The point is that the same minimal data are readable in two equivalent ways: geometrically as an oriented double cover with seam involution, and algebraically as the $3+2$ carrier split with two spectral hypercharges.

\begin{figure}[t]
\centering
\footnotesize
\begin{tikzpicture}[
    >=Latex,
    slot3/.style={draw=blue!60!black, fill=blue!6, rounded corners, minimum width=0.72cm, minimum height=0.72cm, align=center},
    slot2/.style={draw=orange!70!black, fill=orange!9, rounded corners, minimum width=0.72cm, minimum height=0.72cm, align=center},
    info/.style={draw=black!55, rounded corners, fill=gray!3, align=center, inner sep=5pt}
]
\node[font=\bfseries] at (-3.5,1.95) {oriented double cover};
\draw[draw=gray!70!black, fill=gray!6, rounded corners] (-5.4,0.55) rectangle (-1.7,1.35);
\draw[draw=gray!70!black, fill=gray!6, rounded corners] (-5.4,-1.35) rectangle (-1.7,-0.55);
\node at (-3.55,0.95) {upper sheet};
\node at (-3.55,-0.95) {lower sheet};
\draw[very thick, purple!70!black] (-5.6,0) -- (-1.45,0);
\node[anchor=west, purple!70!black] at (-1.38,0) {seam $\Sigma$};
\draw[<->, thick] (-3.55,0.42) -- (-3.55,-0.42);
\node[anchor=west] at (-3.4,0) {$\iota$};

\node[font=\bfseries] at (2.35,1.95) {minimal carrier split};
\node[slot3] (c1) at (0.55,0.95) {$c_1$};
\node[slot3, right=2mm of c1] (c2) {$c_2$};
\node[slot3, right=2mm of c2] (c3) {$c_3$};
\node[slot2, right=6mm of c3] (w1) {$w_1$};
\node[slot2, right=2mm of w1] (w2) {$w_2$};
\node[align=center, blue!60!black] at ($(c2.south)+(0,-0.45)$) {$E_3=P_3E$\\three color slots};
\node[align=center, orange!75!black] at ($(w1.south)+(0.45,-0.45)$) {$E_2=P_2E$\\two weak slots};
\draw[decorate, decoration={brace, amplitude=4pt}] ($(c1.north west)+(-0.07,0.12)$) -- ($(c3.north east)+(0.07,0.12)$);
\node[blue!60!black] at ($(c2.north)+(0,0.42)$) {$3$};
\draw[decorate, decoration={brace, amplitude=4pt}] ($(w1.north west)+(-0.07,0.12)$) -- ($(w2.north east)+(0.07,0.12)$);
\node[orange!75!black] at ($(w1.north east)+(0.2,0.42)$) {$2$};

\node[info] at (-0.45,-2.3) {%
    $\displaystyle
    Y=-\frac13 P_3+\frac12 P_2
    =\frac1{12}\mathbf 1+\frac5{12}\iota
    \qquad
    \iota=\frac{12Y-\mathbf 1}{5}
    $\\[2pt]
    $\displaystyle
    P_3=\frac{\mathbf 1-\iota}{2},
    \qquad
    P_2=\frac{\mathbf 1+\iota}{2},
    \qquad
    \mathrm{Spec}(Y)=\left\{-\frac13,\frac12\right\}
    $};
\end{tikzpicture}
\caption{Minimal carrier at a glance. The seam involution and the hypercharge operator carry the same hard two-point information: geometrically they separate the two sheets, and algebraically they separate the $3+2$ carrier split through the projectors $P_3$ and $P_2$.}
\label{fig:minimal-carrier-glance}
\end{figure}

\begin{remark}[Two-point carrier algebra]
Because the minimal carrier satisfies $6Y^2-Y-\mathbf 1=0$, every analytic function of $Y$ reduces to
\[
f(Y)=A_f\,\mathbf 1+B_f\,Y,
\qquad
A_f=\frac{2f(1/2)+3f(-1/3)}{5},
\qquad
B_f=\frac{6\bigl(f(1/2)-f(-1/3)\bigr)}{5}.
\]
Thus the hard carrier is effectively a two-point algebra supported on the spectral values $-1/3$ and $1/2$, and the projectors $P_3,P_2$ are already polynomial functions of $Y$.
\end{remark}

\begin{remark}[Factorized carrier polynomial and carrier discriminant]\label{rem:carrier-discriminant}
For a general split carrier $E=E_b\oplus E_s$ with total rank $g:=b+s$ and eigenvalues $-1/b$ and $1/s$, the carrier generator satisfies
\[
\left(Y+\frac{1}{b}\right)\left(Y-\frac{1}{s}\right)=0
\quad\Longleftrightarrow\quad
bs\,Y^2+(s-b)Y-\mathbf 1=0.
\]
Its discriminant is
\[
\Delta_Y=(s-b)^2+4bs=(b+s)^2=g^2.
\]
Thus the square root of the carrier discriminant is the carrier rank itself. Moreover,
\[
Y=\frac{b-s}{2bs}\,\mathbf 1+\frac{g}{2bs}\,\inv,
\qquad
\inv=\frac{2bs\,Y-(b-s)\,\mathbf 1}{g}.
\]
On the minimal branch $(b,s)=(3,2)$, this reduces to
\[
6Y^2-Y-\mathbf 1=(2Y-\mathbf 1)(3Y+\mathbf 1),
\qquad
\Delta_Y=25=5^2,
\qquad
\inv=\frac{12Y-\mathbf 1}{5}.
\]
Hence the denominator $5$ in the hypercharge--seam inversion is exactly $\sqrt{\Delta_Y}$, not an additional inserted constant.

This does not replace the stabilizer proof of
\[
G_{\SM}^{\mathrm{glob}}=S(U(3)\times U(2)),
\]
which is the actual source of the global Standard Model quotient in the present manuscript. What the factorization and discriminant do show exactly is that the hard carrier already supports the distinguished spectral values
\[
Y=\frac{1}{2},
\qquad
Y=-\frac{1}{3},
\]
and that the inversion denominator is fixed algebraically by the carrier discriminant.
\end{remark}

The two carrier invariants used repeatedly below are
\[
\gamma=\tr_EY^2=\frac56,
\qquad
B=\frac{\rank E_3}{\rank E_2}=\frac32.
\]
\begin{corollary}[Carrier norm from the carrier polynomial]
The carrier polynomial already fixes
\[
\gamma=\tr_EY^2=\frac16\tr_E(Y+\mathbf 1)=\frac56,
\]
since $\tr_EY=0$ and $\rank E=5$.
\end{corollary}

\begin{corollary}[Minimal carrier compression]
\label{cor:minimal-carrier-compression}
On the minimal carrier branch,
\[
B\gamma=\frac32\cdot\frac56=\frac54.
\]
\end{corollary}

We package the hard kernel as
\[
\Cmin=(E_3\oplus E_2,\ Y,\ S^+,\ \gamma,\ B).
\]

